\begin{mistake}{2026-06-30}{05}

\begin{problemcontent}
设函数 \(f(x)\) 在 \([-l,l]\) 上连续，在 \(x=0\) 处可导，且 \(f'(0)\ne 0\)。

\(1\) 证明：对任意 \(x\in(0,l)\)，至少存在 \(\theta\in(0,1)\)，使得
\[
\int_0^x f(t)\,dt+\int_0^{-x}f(t)\,dt
=x\bigl[f(\theta x)-f(-\theta x)\bigr].
\]

\(2\) 求极限 \(\displaystyle \lim_{x\to0^+}\theta\)。
\end{problemcontent}

\begin{solutioncontent}
令
\[
F(x)=\int_0^x f(t)\,dt+\int_0^{-x}f(t)\,dt.
\]
由于 \(f\) 连续，\(F\) 在相应区间上可导，且
\[
F'(x)=f(x)-f(-x),
\]
其中第二项使用链式法则：
\[
\frac{d}{dx}\int_0^{-x}f(t)\,dt=-f(-x).
\]
又 \(F(0)=0\)。对 \(F\) 在 \([0,x]\) 上用拉格朗日中值定理，存在 \(\xi\in(0,x)\)，使
\[
F(x)-F(0)=xF'(\xi).
\]
令 \(\xi=\theta x\)，则 \(\theta\in(0,1)\)，从而
\[
\int_0^x f(t)\,dt+\int_0^{-x}f(t)\,dt
=x\bigl[f(\theta x)-f(-\theta x)\bigr].
\]

下面求 \(\lim_{x\to0^+}\theta\)。由上式得
\[
\frac{\int_0^x f(t)\,dt+\int_0^{-x}f(t)\,dt}{2x^2}
=
\frac{f(\theta x)-f(-\theta x)}{2\theta x}\cdot \theta.
\]
记
\[
A(x)=\frac{\int_0^x f(t)\,dt+\int_0^{-x}f(t)\,dt}{2x^2},
\qquad
B(x)=\frac{f(\theta x)-f(-\theta x)}{2\theta x}.
\]
则
\[
A(x)=B(x)\theta.
\]

先求 \(A(x)\)。由洛必达法则，
\[
\lim_{x\to0^+}A(x)
=
\lim_{x\to0^+}\frac{f(x)-f(-x)}{4x}.
\]
又
\[
\frac{f(x)-f(-x)}{4x}
=\frac14\left[\frac{f(x)-f(0)}{x}+\frac{f(0)-f(-x)}{x}\right].
\]
因为 \(f\) 在 \(0\) 处可导，所以
\[
\lim_{x\to0^+}A(x)=\frac14\bigl[f'(0)+f'(0)\bigr]=\frac12f'(0).
\]

再求 \(B(x)\)。令 \(y=\theta x\)。由于 \(0<\theta<1\)，当 \(x\to0^+\) 时，\(y\to0^+\)。因此
\[
\lim_{x\to0^+}B(x)
=
\lim_{y\to0^+}\frac{f(y)-f(-y)}{2y}.
\]
并且
\[
\frac{f(y)-f(-y)}{2y}
=\frac12\left[\frac{f(y)-f(0)}{y}+\frac{f(0)-f(-y)}{y}\right]
\to f'(0).
\]
故
\[
\lim_{x\to0^+}B(x)=f'(0)\ne0.
\]
由 \(A(x)=B(x)\theta\)，可得
\[
\theta=\frac{A(x)}{B(x)}.
\]
于是
\[
\lim_{x\to0^+}\theta
=\frac{\frac12f'(0)}{f'(0)}
=\frac12.
\]
因此
\[
\boxed{\lim_{x\to0^+}\theta=\frac12}.
\]
\end{solutioncontent}

\mistakeknowledge{
\begin{itemize}
  \item \textbf{核心公式/定理}：拉格朗日中值定理：\(F(x)-F(0)=xF'(\xi)\)，其中 \(\xi\in(0,x)\)；变上限积分求导需结合链式法则。
  \item \textbf{易错点}：不能在尚未证明 \(\lim\theta\) 存在时直接使用乘积极限定理写成 \(\lim(B\theta)=\lim B\cdot\lim\theta\)。严谨做法是先证 \(A(x)\to \frac12f'(0)\)、\(B(x)\to f'(0)\ne0\)，再由 \(\theta=A/B\) 得极限。
  \item \textbf{关联知识}：若 \(f\) 在 \(0\) 处可导，则 \(\displaystyle \lim_{y\to0^+}\frac{f(y)-f(-y)}{2y}=f'(0)\)。
\end{itemize}}

\end{mistake}
